\section{The \texttt{Pol\_Bfield\_stop} McStas Component}
Magnetic field component.

\subsection*{Identification}
\begin{itemize}
  \item \textbf{Site:} 
  \item \textbf{Author:} Erik B Knudsen, Peter Christiansen, and Peter Willendrup
  \item \textbf{Origin:} RISOE
  \item \textbf{Date:} August 2006
\end{itemize}

\subsection*{Description}
\begin{lstlisting}
End of magnetic field region defined by the latest preceeding Pol_Bfield component.

The component is concentric. It means that it requires a

// START MAGNETIC FIELD
COMPONENT msf =
Pol_Bfield(xw=0.08, yh=0.08, length=0.2, Bx=0, By=-0.678332e-4, Bz=0)
AT (0, 0, 0) RELATIVE armMSF

// HERE CAN BE OTHER COMPONENTS INSIDE THE MAGNETIC FIELD

// STOP MAGNETIC FIELD
COMPONENT msfCp = Pol_Bfield_stop()
AT ("SOMEWHERE") RELATIVE armMSF

In between the two components the propagation routine
PROP_DT also handles the spin propagation.
The current algorithm used for spin propagation is:
SimpleNumMagnetPrecession
in pol-lib.
and does not handle gravity.

GRAVITY: NO
POLARISATION: YES

Example: Pol_Bfield_stop()
\end{lstlisting}

\subsection*{Input parameters}
Parameters in \textbf{boldface} are required; the others are optional.

\begin{longtable}{p{0.22\textwidth}p{0.12\textwidth}p{0.46\textwidth}p{0.14\textwidth}}
\toprule
\textbf{Name} & \textbf{Unit} & \textbf{Description} & \textbf{Default} \\
\midrule
\endhead
geometry & str & Name of an Object File Format (OFF) or PLY file for complex field-geometry. & "" \\
yheight & m & Height of opening. & 0 \\
xwidth & m & Width of opening. & 0 \\
zdepth & m & Length of field. & 0 \\
radius & m & Radius of field if it is cylindrical or spherical. & 0 \\
\bottomrule
\end{longtable}

\subsection*{Links}
\begin{itemize}
  \item \href{run:/home/willend/willend-McCode/mcstas-comps/optics/Pol_Bfield_stop.comp}{Source code} for \texttt{Pol\_Bfield\_stop.comp}.
\end{itemize}
\IfFileExists{Pol_Bfield_stop_static.tex}{\input{Pol_Bfield_stop_static.tex}}{}